% Research contribution: number fields.  This is an input fragment for the
% principal report; theorem/proposition/lemma/remark environments are assumed.

\section{Number fields: what transfers, and what does not}
\label{sec:number-fields}

Let $K$ be a number field, that is, a finite extension of $\mathbb Q$, and let
$\mathcal O_K$ be its ring of algebraic integers.  Coefficients in $K$ are
allowed in a Diophantine definition over $K$.  This section distinguishes
the existence of a definition for each $K$ from the much stronger assertion
that one fixed equation works after every extension of scalars.

\subsection{An exact counterexample to extending the released polynomial}

Write $P(X_0,\ldots,X_{12503})$ for the released quartic, with the variable
indices used in its monomial table, and put
\[
 C=5937913971887799417367.
\]
The following calculation is independent of the arithmetic argument that
$P$ defines $\mathbb Z$ over $\mathbb Q$.

\begin{proposition}\label{prop:quartic-gaussian}
Over $\mathbb Q(i)$, the existential projection of the released polynomial
$P$ onto its input coordinate is all of $\mathbb Q(i)$.  In particular,
this polynomial does not define $\mathbb Z[i]$ in $\mathbb Q(i)$.
\end{proposition}
\begin{proof}
Set every witness variable other than $X_{2242}=u$ and $X_{2243}=v$ equal
to zero.  The exact resulting polynomial is
\begin{equation}\label{eq:gaussian-restriction}
 P(t,0,\ldots,u,\ldots,v,\ldots,0)=C+4t^2+t^4+u^2+v^2.
\end{equation}
For any $t$ in a field containing $\mathbb Q(i)$, put
\[
 R=C+4t^2+t^4,\qquad u=\frac{1-R}{2},\qquad
 v=i\frac{1+R}{2}.
\]
Then
\[
 R+u^2+v^2
 =R+\frac{(1-R)^2-(1+R)^2}{4}=0.
\]
Thus every $t\in\mathbb Q(i)$ is accepted.  For example, $t=1/2$ is not
in $\mathbb Z[i]$ and has the witnesses
\[
 X_{2242}=-\frac{95006623550204790677873}{32},\qquad
 X_{2243}=i\frac{95006623550204790677905}{32},
\]
with all other witness coordinates zero.

For reproducibility, \texttt{check\_base\_change.py} checks
\eqref{eq:gaussian-restriction} in two independent ways: by retaining the
relevant monomials of the expanded table, and by specializing and squaring
the quadratic residuals.  It also substitutes the displayed generic
witnesses directly in both representations, using exact arithmetic in
$\mathbb Q(i)[t]$.  The source table has SHA-256
\begin{center}\small\ttfamily
7f5aa3117af82a0b323f0a4e74a18599204f5540408106d3ce99612b5ed830ca.
\end{center}
\end{proof}

\begin{remark}
Proposition~\ref{prop:quartic-gaussian} is a counterexample to direct
extension of the particular equation.  It is not a counterexample to
the existence of a different Diophantine definition of
$\mathbb Z[i]$ in $\mathbb Q(i)$.
The obstruction is already present in the final sum-of-squares operation:
$u^2+v^2=0$ no longer forces $u=v=0$ over $\mathbb Q(i)$.
Even replacing this operation by a suitable norm form would not, by
itself, establish that the original arithmetic system characterizes
integrality over a new field.
\end{remark}

\subsection{A general obstruction to uniform scalar extension}

\begin{proposition}\label{prop:bounded-extension-false-positive}
Suppose that $P\in\mathbb Q[T,W_1,\ldots,W_m]$ has total degree $D$ and
defines $\mathbb Z$ in $\mathbb Q$.  There is a finite set
$E\subset\mathbb Q$ with $|E|\leq D-1$ such that, for every
$q\in\mathbb Q\setminus(\mathbb Z\cup E)$, there exist a number field
$K_q$ and $\mathbf w\in K_q^m$ satisfying
\[
 2\leq[K_q:\mathbb Q]\leq D,\qquad P(q,\mathbf w)=0,
 \qquad q\notin\mathcal O_{K_q}.
\]
Consequently no fixed polynomial defining $\mathbb Z$ in $\mathbb Q$
defines $\mathcal O_K$ in every number field $K$ by scalar extension.
\end{proposition}
\begin{proof}
The polynomial must depend on at least one $W_j$: otherwise its rational
zero set is finite or all of $\mathbb Q$.  Expand it as
\[
 P(T,\mathbf W)=\sum_{\alpha}c_{\alpha}(T)\mathbf W^{\alpha}
\]
and select $\alpha\ne0$ with $c_{\alpha}\ne0$.  Let $E$ be the rational
zero set of $c_{\alpha}$.  Since $\deg c_{\alpha}\leq D-|\alpha|$,
we have $|E|\leq D-1$.  If $q\notin E$, the polynomial
$P(q,\mathbf W)$ is nonconstant.

Choose a variable on which this polynomial depends.  Its leading
coefficient in that variable is a nonzero polynomial in the remaining
variables.  A nonzero polynomial over an infinite field cannot vanish
at every rational point; by induction on the number of variables,
we can therefore specialize all the remaining variables to rational
numbers so that the resulting polynomial $h(U)\in\mathbb Q[U]$ is
nonconstant.  Its degree is at most $D$.
Choose an irreducible factor of $h$, let $\theta$ be a root, and set
$K_q=\mathbb Q(\theta)$.  The specialized coordinates, with $U=\theta$,
give the required zero of $P$.  The degree of $K_q$ is at most $D$.
It cannot be one because $q\notin\mathbb Z$ and $P$ defines
$\mathbb Z$ over $\mathbb Q$.

Finally, $\mathcal O_{K_q}\cap\mathbb Q=\mathbb Z$: a rational
algebraic integer has denominator one, by the rational-root theorem
applied to a monic integral equation.  Hence $q$ is not integral.
\end{proof}

\begin{remark}
If one chosen witness has degree at most $d$ in $P$ and genuinely occurs
after specialization, the same proof gives $[K_q:\mathbb Q]\leq d$.
For a sum of squares whose residuals are affine linear in that witness,
a nonconstant quadratic specialization with no rational zero has
strictly negative discriminant.  Indeed it is a sum of squares of real
affine linear polynomials, so its discriminant is nonpositive, and
discriminant zero would give a rational root.  In this case an imaginary
quadratic extension already supplies a false positive.
\end{remark}

\subsection{Combining equations over an arbitrary number field}

The use of one polynomial rather than a finite system causes no
definability obstruction over a number field, although it can affect
the degree severely.

\begin{lemma}\label{lem:norm-combine}
Let $K$ be a number field, let $s\geq1$, and let
$f_1,\ldots,f_s\in K[\mathbf X]$ have degree at most $d$.
There is a polynomial $F\in K[\mathbf X]$, with no additional variables
and degree at most $sd$, such that
\[
 F(\mathbf x)=0\quad\Longleftrightarrow\quad
 f_1(\mathbf x)=\cdots=f_s(\mathbf x)=0
 \qquad(\mathbf x\in K^n).
\]
\end{lemma}
\begin{proof}
Choose a field extension $L/K$ of degree $s$ and a $K$-basis
$\beta_1,\ldots,\beta_s$ of $L$.  Such an extension exists: choose a
discrete valuation of $K$ with a uniformizer $\pi\in K$; the polynomial
$U^s-\pi$ is Eisenstein.  The norm form
\[
 N(Z_1,\ldots,Z_s)=
 N_{L/K}\!\left(\sum_{j=1}^s\beta_j Z_j\right)
\]
is homogeneous of degree $s$ over $K$, and it vanishes on a vector in
$K^s$ precisely when that vector is zero.  Take
$F=N(f_1,\ldots,f_s)$.
\end{proof}

For the degree objective, the preceding elementary norm bound can be
substantially improved.  The following construction uses only a single
nonarchimedean place.

\begin{lemma}\label{lem:four-variable-anisotropic}
Every number field $K$ admits an anisotropic quadratic form in four
variables with coefficients in $\mathcal O_K$.  Consequently a system
of $s\geq2$ polynomial equations of degree at most $d$ over $K$ can be
combined, without additional variables, into one equation of degree at
most
\[
 2^{\lceil\log_4 s\rceil}d.
\]
If $K$ admits a real embedding, degree at most $2d$ suffices for every
$s$.
\end{lemma}
\begin{proof}
Choose a prime ideal $\mathfrak p$ of $\mathcal O_K$ above an odd
rational prime, and let $v$ be its normalized discrete valuation.
Choose $a\in\mathcal O_K$ whose image in
$\mathcal O_K/\mathfrak p$ is a nonsquare, and choose
$\pi\in\mathfrak p\setminus\mathfrak p^2$.  For any
$(x,y)\ne(0,0)$ in $K^2$, reduction modulo $\mathfrak p$ after
division by $\pi^{\min\{v(x),v(y)\}}$ gives
\[
 v(x^2-ay^2)=2\min\{v(x),v(y)\}.
\]
Indeed the two leading residues cannot cancel because $a$ is a
nonsquare in the residue field.  It follows that
\[
 q(X_1,X_2,X_3,X_4)
 =X_1^2-aX_2^2-\pi X_3^2+a\pi X_4^2
\]
is anisotropic: equality of its two binary parts would equate an even
valuation with an odd one, unless both parts vanish, in which case all
four coordinates are zero.

Place the $s$ residuals at the leaves of a complete tree of branching
degree four and height $h=\lceil\log_4 s\rceil$, padding the unused
leaves by zero.  At every internal vertex apply $q$ to its four child
polynomials.  Anisotropy shows inductively that the resulting polynomial
vanishes exactly when every original residual vanishes.  Each level
at most doubles the degree, giving the stated bound.  If $K$ has a
real embedding, the sum of the squares of all the residuals works
directly, since a sum of real squares is zero only when every summand
is zero.
\end{proof}

\subsection{Precise transfer principles}

\begin{lemma}[Transitivity]\label{lem:diophantine-transitivity-nf}
Let $A\subset B\subset C$ be rings of characteristic zero whose
fraction fields are not algebraically closed.  If $A$ is Diophantine in
$B$ and $B$ is Diophantine in $C$, then $A$ is Diophantine in $C$.
Here the coefficients used to define $A$ are regarded as constants in $C$.
\end{lemma}
\begin{proof}
Take a defining existential system for $A$ over $B$.  Regard its
variables as ranging over $C$, and append a defining system for
membership in $B$ to every variable, including the free variable.
The resulting finite existential system defines exactly $A$ in $C$.
It can be combined into one equation by an anisotropic norm form,
as in Lemma~\ref{lem:norm-combine} for the number-field cases used here.
For the stated more general rings, the same construction uses the
homogenization of an irreducible polynomial over their fraction fields
and clears denominators.
\end{proof}

\begin{proposition}\label{prop:integral-basis-transfer}
If $\mathbb Z$ is Diophantine in a number field $K$ of degree $n$, then
$\mathcal O_K$ is Diophantine in $K$.
\end{proposition}
\begin{proof}
Choose an integral basis $\omega_1,\ldots,\omega_n$ of $\mathcal O_K$.
Then
\[
 x\in\mathcal O_K\quad\Longleftrightarrow\quad
 \exists z_1,\ldots,z_n\in\mathbb Z:\quad
 x=\sum_{j=1}^n\omega_jz_j.
\]
Replace each condition $z_j\in\mathbb Z$ by its Diophantine
definition over $K$, and combine the finitely many equations by
Lemma~\ref{lem:norm-combine}.
\end{proof}

\begin{proposition}[Quantitative integral-basis construction]
\label{prop:quantitative-integral-basis}
Suppose a polynomial $G(T,\mathbf W)$ of degree $d$ defines
$\mathbb Z$ over $K$, using $m$ witnesses.  If $[K:\mathbb Q]=n$,
then $\mathcal O_K$ has a definition using at most
\[
 n(m+1)-1
\]
witnesses.  For $n\geq2$ the total degree can be at most
$2^{\lceil\log_4 n\rceil}d$, and at most $2d$ when $K$ admits a
real embedding.  For $n=1$, the original polynomial suffices.
\end{proposition}
\begin{proof}
In the integral-basis construction, eliminate the last coefficient by
\[
 z_n=\frac{x-\sum_{j<n}\omega_j z_j}{\omega_n}.
\]
Apply $G$ to each of $z_1,\ldots,z_n$, using separate witness tuples.
The remaining coefficient variables number $n-1$ and the witnesses
inside the $n$ copies of $G$ number $nm$.  Substitution of the displayed
linear expression for $z_n$ does not increase degree, and multiplication
by a nonzero element of $K$ clears any coefficient denominator.  There
are now $n$ equations of degree at most $d$; apply
Lemma~\ref{lem:four-variable-anisotropic}.
\end{proof}

\begin{proposition}\label{prop:descent-integrality}
Let $L/K$ be an extension of number fields.  If $\mathcal O_L$ is
Diophantine in $L$, then $\mathcal O_K$ is Diophantine in $K$.
\end{proposition}
\begin{proof}
Choose a $K$-basis $\beta_1=1,\beta_2,\ldots,\beta_e$ of $L$.
In a defining polynomial for $\mathcal O_L$, replace each quantified
$L$-variable by $\sum_{j=1}^e\beta_j y_j$, with the $y_j$ quantified
over $K$.  Expand every coefficient in the same basis and equate the
$e$ coordinates of the resulting expression to zero.  For a free
variable $x\in K$, this system is equivalent to $x\in\mathcal O_L$,
which is equivalent to $x\in\mathcal O_K$.  Combine the equations
using Lemma~\ref{lem:norm-combine}.
\end{proof}

\subsection{The remaining arithmetic question}

The theorem of Koymans--Pagano and, independently,
Alp\"oge--Bhargava--Ho--Shnidman proves that
\begin{equation}\label{eq:dl-theorem}
 \mathbb Z\text{ is Diophantine in }\mathcal O_K
 \quad\text{for every number field }K.
\end{equation}
This is Theorem~1.2 of \cite{KoymansPagano2026} and Corollary~1.2 of
\cite{AlpogeBhargavaHoShnidman2026}.  These statements concern witnesses
in $\mathcal O_K$; they do not supply unrestricted witnesses in $K$.

\begin{theorem}[Conditional equivalence of the extension questions]
\label{thm:nf-equivalence}
Assume that $\mathbb Z$ is Diophantine in $\mathbb Q$.
For a number field $K$, the following are equivalent:
\begin{enumerate}
 \item $\mathcal O_K$ is Diophantine in $K$;
 \item $\mathbb Z$ is Diophantine in $K$;
 \item $\mathbb Q$ is Diophantine in $K$.
\end{enumerate}
The equivalence of the first two assertions uses the established
theorem~\eqref{eq:dl-theorem} but does not require the assumption about
$\mathbb Q$.
\end{theorem}
\begin{proof}
The first assertion implies the second by
\eqref{eq:dl-theorem} and transitivity.  The second implies the first
by Proposition~\ref{prop:integral-basis-transfer}.
For the implication from the second to the third, use
\[
 x\in\mathbb Q\quad\Longleftrightarrow\quad
 \exists a,b\in\mathbb Z\ \exists c\in K:
 \quad bx=a,\quad bc=1.
\]
Finally, if $\mathbb Q$ is Diophantine in $K$, transitivity transports
the assumed definition of $\mathbb Z$ over $\mathbb Q$ to $K$.
\end{proof}

These transfer statements do not construct a Diophantine definition of
$\mathbb Q$ in an arbitrary number field.  The rank-stability theorems
behind~\eqref{eq:dl-theorem} establish definitions over rings of integers,
not that missing field-level assertion.  The results proved here give
no counterexample to the existence of a field-dependent definition of
$\mathcal O_K$, and do not settle that general question.  They do show
that an extension argument must contain new arithmetic input: neither
the reuse of the released equation nor an unrestricted change of the
witness field suffices.
